\vssub
\subsection{~传播}
\vssub

在数值模型中，需要采用平衡方程（\ref{eq:balance0}）的 Euler 形式。该平衡
方程既可以写成输运方程形式（速度位于导数之外），也可以写成守恒形式
（速度位于导数之内）。前一种形式只适用于向量波数谱 $N(\bk;\bx,t)$；后一种
形式则可针对任意谱表达推导出有效方程，只要前述相应的 Jacobian 变换表现
良好即可 \citep[例如][]{tol:GAOS98b}。此外，与输运方程不同，守恒方程保持
总波能/作用量守恒。当方程用于数值模型时，这是一个重要特性。\ws\ 中使用的
谱 $N(k,\theta;\bx,t)$ 的平衡方程为（为简化记号，下文将该谱简写为 $N$）：

%----------------------------------%
% Plain-grid propagation equations %
%----------------------------------%
% eq:bal_plane
% eq:x_dot
% eq:k_dot
% eq:theta_dot

\begin{equation}
\frac{\partial N}{\partial t} + \nabla_x \cdot {\bf\dot{x}}N +
\frac{\partial}{\partial k} \dot{k}N +
\frac{\partial}{\partial \theta} \dot{\theta}N =
\frac{S}{\sigma} \: , \label{eq:bal_plane}
\end{equation}
\begin{equation}
{\bf \dot{x}} = {\bf c}_g + {\bf U} \: , \label{eq:x_dot}
\end{equation}
\begin{equation}
\dot{k} = - \frac{\partial \sigma}{\partial d}
\frac{\partial d}{\partial s} - {\bk} \cdot
\frac{\partial {\bf U}}{\partial s} \: , \label{eq:k_dot}
\end{equation}
\begin{equation}
\dot{\theta} = - \frac{1}{k} \left [
\frac{\partial \sigma}{\partial d} \frac{\partial d}{\partial m}
+ {\bk} \cdot \frac{\partial {\bf U}}{\partial m}
\right ] \: , \label{eq:theta_dot}
\end{equation}

\noindent
其中 ${\bf c}_g = (c_g \sin \theta, c_g \cos \theta)$，$s$ 是沿
$\theta$ 方向的坐标，$m$ 是垂直于 $s$ 的坐标。式~(\ref{eq:bal_plane})
适用于 Cartesian 坐标。对于大尺度应用，该方程通常转换为由经度 $\lambda$
和纬度 $\phi$ 定义的 spherical 坐标，同时保持局地方差的定义
\citep[即单位表面积上的方差，如][]{art:WAM88}：

%---------------------------------%
% Spherical propagation equations %
%---------------------------------%
% eq:bal_sphere
% eq:phi_dot
% eq:lambda_dot
% eq:theta_g_dot

\begin{equation}
\frac{\partial N}{\partial t} +
\frac{1}{\cos \phi} \frac{\partial}{\partial \phi} \dot{\phi}
    N \cos \theta +
\frac{\partial}{\partial \lambda} \dot{\lambda}N +
\frac{\partial}{\partial k} \dot{k} N +
\frac{\partial}{\partial \theta} \dot{\theta}_g N
= \frac{S}{\sigma} \: ,\label{eq:bal_sphere}
\end{equation}
\begin{equation}
\dot{\phi} = \frac{c_g \cos \theta + U_\phi}{R}
\: , \label{eq:phi_dot}
\end{equation}
\begin{equation}
\dot{\lambda} =  \frac{c_g \sin \theta + U_\lambda}{R \cos \phi}
\: , \label{eq:lambda_dot}
\end{equation}
\begin{equation}
\dot{\theta}_g = \dot{\theta} -
\frac{c_g \tan \phi \cos \theta}{R}
\: , \label{eq:theta_g_dot}
\end{equation}

\noindent
其中 $R$ 是地球半径，$U_\phi$ 和 $U_\lambda$ 是流速分量。
式~(\ref{eq:theta_g_dot}) 包含沿大圆传播的修正项，并使用 $\theta$ 的
Cartesian 定义；其中 $\theta = 0$ 对应从西向东传播的波。\ws\ 可使用
Cartesian 或 spherical 坐标运行。注意，岛屿等未解析障碍物也可纳入方程。
在 \ws\ 中，这是在数值格式层面完成的，见第~\ref{sub:num_obst} 节。
同样，波长尺度上的水深变化可通过第~\ref{sec:BS1} 节所述的散射源项引入。

最后，Cartesian 坐标和 spherical 坐标都可以用多种方式离散，包括四边形
（矩形网格、曲线网格或 SMC 网格）和三角形。这一部分在
第~\ref{chapt:num} 章中讨论。
